\documentclass[11pt,a4paper]{article}

\usepackage[T1]{fontenc}
\usepackage[utf8]{inputenc}
\usepackage[french]{babel}
\usepackage{lmodern}
\usepackage[a4paper,margin=2.1cm]{geometry}
\usepackage{setspace}
\usepackage{hyperref}
\usepackage{booktabs}
\usepackage{longtable}
\usepackage{array}
\usepackage{enumitem}
\usepackage{xcolor}
\usepackage{tikz}
\usepackage{titlesec}
\usepackage{fancyhdr}

\usetikzlibrary{positioning,arrows.meta,shapes.geometric,fit,backgrounds}

\hypersetup{
  colorlinks=true,
  linkcolor=blue!50!black,
  urlcolor=blue!50!black,
  pdftitle={Rapport Projet Containerisation Web},
  pdfauthor={Rayane CHERIFI \& SEBA Abdelwahab}
}

\setstretch{1.08}
\setlist[itemize]{leftmargin=1.4em,itemsep=0.2em,topsep=0.25em}
\setlist[enumerate]{leftmargin=1.6em,itemsep=0.2em,topsep=0.25em}

\pagestyle{fancy}
\fancyhf{}
\lhead{Projet Containerisation Web}
\rhead{Movie Catalog}
\cfoot{\thepage}

\titleformat{\section}{\large\bfseries}{\thesection.}{0.5em}{}
\titleformat{\subsection}{\normalsize\bfseries}{\thesubsection.}{0.5em}{}

\begin{document}

\begin{titlepage}
  \centering
  {\Large \textbf{Rapport de Projet}}\\[0.4cm]
  {\LARGE \textbf{Conception et Deploiement d'une Application Web}}\\[0.2cm]
  {\LARGE \textbf{Orientee Donnees et API}}\\[1.2cm]

  {\large UE: CRV - Containerisation et Deploiement}\\[0.3cm]
  {\large Annee universitaire 2025-2026}\\[1.0cm]

  \begin{tabular}{rl}
    \textbf{Projet:} & Movie Catalog Web Application \\
    \textbf{Binome:} & Rayane CHERIFI, SEBA Abdelwahab \\
    \textbf{Date limite:} & 17 mai 2026 \\
    \textbf{Date du rapport:} & \today \\
  \end{tabular}

  \vfill

  \begin{minipage}{0.95\textwidth}
  \small
  Ce rapport presente la conception, l'implementation et le deploiement d'une application web complete,
  conformement au cahier des charges du sujet \textit{``Conception et deploiement d'une application Web orientee donnees et API''}.
  L'accent est mis sur la justification technique, la robustesse du deploiement et la reproductibilite du lancement.
  \end{minipage}
\end{titlepage}

\tableofcontents
\newpage

\section{Contexte, objectifs et perimetre}

Le projet s'inscrit dans une logique professionnelle: il ne s'agit pas uniquement de coder une interface,
mais de livrer un systeme deployable de bout en bout (front-end, back-end, donnees, conteneurisation, orchestration,
validation, documentation).

\subsection{Objectif fonctionnel}

L'application \textbf{Movie Catalog} permet:
\begin{itemize}
  \item d'afficher un catalogue de films issu d'un dataset reel (TMDB 5000),
  \item de rechercher des films (recherche simple + recherche avancee),
  \item de consulter les details d'un film,
  \item de gerer des favoris pour des utilisateurs authentifies,
  \item d'exposer ces donnees via une API HTTP structuree.
\end{itemize}

\subsection{Objectif technique}

Le projet devait absolument demontrer:
\begin{itemize}
  \item une architecture claire \textit{front-end / back-end / base de donnees},
  \item l'usage concret de Docker et Kubernetes,
  \item la gestion de configuration (ConfigMaps/Secrets),
  \item une demarche d'automatisation (scripts de lancement + CI/CD),
  \item une capacite de justification des choix techniques.
\end{itemize}

\section{Conformite explicite au cahier des charges}

Le tableau suivant mappe les exigences du sujet avec des preuves concretes d'implementation.

\begin{longtable}{p{0.33\textwidth} p{0.61\textwidth}}
\toprule
\textbf{Exigence du sujet} & \textbf{Reponse du projet} \\
\midrule
\endhead
Catalogue de films dynamique & Endpoint REST \texttt{/api/movies/} avec pagination DRF, affichage React (grille de films). \\
Barre de recherche & Recherche standard (parametre \texttt{search}) + recherche avancee via Elasticsearch (endpoint \texttt{/api/search/}). \\
Architecture front/back/data & React + Vite (front), Django REST (back), PostgreSQL + Redis + Elasticsearch (data/cache/recherche). \\
API HTTP exposee et consommee & API DRF documentee, consommee via Axios (intercepteurs JWT, refresh automatique des tokens). \\
Conteneurisation Docker & Dockerfiles multi-stage front/back + \texttt{docker-compose.yml} pour environnement local. \\
Deploiement Kubernetes & Kustomize overlay local, Deployments, StatefulSets, Services, Ingress, ConfigMap, Secret, NetworkPolicies. \\
CI/CD & Workflow GitHub Actions: checks backend/frontend, validation K8s render, build d'images Docker, publication manuelle. \\
Script de lancement simple & Scripts \texttt{run/validate/stop} pour Windows, Linux, macOS avec verification de sante. \\
\bottomrule
\end{longtable}

\section{Architecture technique proposee}

\subsection{Schema d'architecture (professionnel et complet)}

\begin{figure}[h!]
\centering
\begin{tikzpicture}[
  node distance=0.9cm and 1.1cm,
  box/.style={draw, rounded corners, align=center, minimum width=2.8cm, minimum height=1cm, fill=blue!5},
  infra/.style={draw, rounded corners, align=center, minimum width=3.0cm, minimum height=1cm, fill=green!6},
  data/.style={draw, rounded corners, align=center, minimum width=2.7cm, minimum height=1cm, fill=orange!8},
  sec/.style={draw, rounded corners, align=center, fill=red!5},
  arr/.style={-Latex, thick}
]

\node[box] (client) {Navigateur\\Utilisateur};
\node[infra, below=of client] (ing) {Ingress NGINX\\(host: movie.local)};
\node[box, below left=of ing] (front) {Frontend\\React + Vite\\(Nginx)};
\node[box, below right=of ing] (back) {Backend API\\Django REST\\Gunicorn};

\node[data, below left=1.1cm and 0.1cm of back] (pg) {PostgreSQL 15\\StatefulSet + PVC};
\node[data, below=1.1cm of back] (redis) {Redis 7\\cache/session};
\node[data, below right=1.1cm and 0.1cm of back] (es) {Elasticsearch 8.11\\index recherche};

\node[infra, right=2.0cm of back] (boot) {Jobs Kubernetes\\migration + bootstrap\\import + index};
\node[sec, below=0.6cm of es] (netpol) {NetworkPolicies\\default deny + flux explicites};
\node[sec, above=0.6cm of pg] (conf) {ConfigMap + Secrets\\(config et credentials)};

\draw[arr] (client) -- (ing);
\draw[arr] (ing) -- node[left,sloped,above]{\small /} (front);
\draw[arr] (ing) -- node[right,sloped,above]{\small /api} (back);
\draw[arr] (front) -- node[above]{\small HTTP/JSON} (back);

\draw[arr] (back) -- (pg);
\draw[arr] (back) -- (redis);
\draw[arr] (back) -- (es);

\draw[arr,dashed] (boot) -- (back);
\draw[arr,dashed] (boot) -- (pg);
\draw[arr,dashed] (boot) -- (es);

\begin{scope}[on background layer]
  \node[draw=gray!50, rounded corners, fit=(ing) (front) (back) (pg) (redis) (es) (boot) (netpol) (conf), inner sep=0.35cm] (cluster) {};
  \node[above right=0.05cm and -0.2cm of cluster.north west, anchor=south west] {\small Cluster Kubernetes (namespace \texttt{movie-catalog})};
\end{scope}
\end{tikzpicture}
\caption{Architecture logique et deploiement cible}
\end{figure}

\subsection{Lecture architecturale}

L'architecture est volontairement \textbf{separee par responsabilite}:
\begin{itemize}
  \item le front gere l'experience utilisateur et la presentation,
  \item le back centralise les regles metier et l'exposition API,
  \item PostgreSQL garantit la persistence relationnelle,
  \item Redis optimise les temps d'acces,
  \item Elasticsearch apporte une recherche plus riche que le simple \texttt{ILIKE} SQL,
  \item Kubernetes orchestre la disponibilite, la resilience et l'isolation reseau.
\end{itemize}

Ce decoupage est realiste en contexte professionnel et rend chaque composant testable/maintenable independamment.

\section{Justification concrete des choix technologiques}

\subsection{Back-end: Django REST Framework}

\textbf{Arguments concrets:}
\begin{itemize}
  \item DRF fournit rapidement des endpoints robustes (ViewSets, serialisation, pagination standard).
  \item La securite native Django est exploitee (validation, middleware de securite, gestion utilisateurs).
  \item Le code expose des endpoints metier clairs: \texttt{/api/movies/}, \texttt{/api/favorites/}, \texttt{/api/health/}, \texttt{/api/search/}.
  \item L'application integre JWT (access+refresh) avec refresh automatique cote front.
\end{itemize}

\textbf{Compromis assume:} Django est plus verbeux que FastAPI pour certains cas, mais il reduit le risque de dette technique
sur un projet avec authentification, administration et persistance relationnelle.

\subsection{Front-end: React + Vite}

\textbf{Arguments concrets:}
\begin{itemize}
  \item React permet la reutilisation de composants (cards, modal, pagination, favoris).
  \item Vite accelere le cycle de developpement et le build de production.
  \item L'API est consommee via Axios avec gestion centralisee des erreurs et des tokens.
  \item Le front implemente des etats utiles a la demo: banniere de sante API, parcours login/register/verification.
\end{itemize}

\subsection{Donnees: PostgreSQL + Redis + Elasticsearch}

\textbf{Arguments concrets:}
\begin{itemize}
  \item \textbf{PostgreSQL} correspond au besoin relationnel (films, utilisateurs, favoris) avec index metier.
  \item \textbf{Redis} est utilise comme cache applicatif (liste/detail/requetes frequentes), ce qui baisse la charge SQL.
  \item \textbf{Elasticsearch} est reserve a la recherche avancee (filtres genres/annee/note, pagination, autocomplete).
\end{itemize}

\textbf{Impact pedagogique et professionnel:} la combinaison SQL + cache + moteur de recherche reproduit un pattern d'architecture
reel pour applications data-centric.

\subsection{Orchestration: Kubernetes}

\textbf{Arguments concrets observes dans les manifests:}
\begin{itemize}
  \item \texttt{Deployment} front/back en \textbf{2 replicas} chacun (rolling update, disponibilite).
  \item \texttt{StatefulSet} pour PostgreSQL et Elasticsearch (identite stable + volumes persistants).
  \item \texttt{ConfigMap} pour la configuration non sensible, \texttt{Secret} pour mots de passe et cles.
  \item \texttt{Ingress} pour le routage \texttt{/} vers front et \texttt{/api} vers back.
  \item \texttt{NetworkPolicy} en mode \textit{default deny} puis ouvertures minimales explicites.
\end{itemize}

Ce choix est coherent avec le sujet qui impose Docker + Kubernetes et valorise la justification de l'architecture.

\section{Fonctionnalites realisees et demonstration}

\subsection{Fonctionnalites minimales (obligatoires)}

\begin{itemize}
  \item Catalogue dynamique des films (API paginee + affichage front).
  \item Barre de recherche (parametre \texttt{search} + page dediee de recherche avancee).
  \item Communication front/back fonctionnelle (health-check affiche dans l'UI).
  \item Gestion d'erreurs minimale (retours API codes HTTP, erreurs front capturees).
\end{itemize}

\subsection{Fonctionnalites avancees (valorisees)}

\begin{itemize}
  \item Favoris utilisateur (creation/suppression, protection par authentification).
  \item Recherche avancee et autocomplete avec Elasticsearch.
  \item Cache Redis sur endpoints critiques.
  \item Pagination configurable et tri cote API.
  \item Authentification JWT + verification email.
\end{itemize}

\subsection{Source de donnees et import}

Le projet exploite \texttt{tmdb\_5000\_movies.csv}. L'import est automatise via commande Django,
puis complete par un indexage Elasticsearch. En environnement Kubernetes, un job \texttt{movies-bootstrap}:
\begin{enumerate}
  \item attend la disponibilite DB,
  \item importe les films si la table est vide,
  \item indexe vers Elasticsearch avec retries,
  \item echoue explicitement si l'indexation n'aboutit pas.
\end{enumerate}

Ce comportement rend le systeme relancable sans reconfiguration manuelle, ce qui est un point fort pour l'evaluation.

\section{Conteneurisation, deploiement et exploitation}

\subsection{Docker}

\begin{itemize}
  \item Dockerfile backend multi-stage (builder Python puis runtime minimal), execution non-root.
  \item Dockerfile frontend multi-stage (build Node, runtime Nginx + healthcheck).
  \item \texttt{docker-compose.yml} pour le mode local (db, redis, elasticsearch, backend, frontend).
\end{itemize}

\subsection{Kubernetes et robustesse operationnelle}

Le deploiement Kubernetes inclut des mecanismes de fiabilite concrets:
\begin{itemize}
  \item probes \texttt{startup/readiness/liveness} sur front et back,
  \item retries de jobs (migration, bootstrap),
  \item \texttt{PodDisruptionBudget} pour front/back,
  \item \texttt{CronJob} de backup PostgreSQL toutes les 6 heures,
  \item ressources requests/limits definies pour chaque composant.
\end{itemize}

\subsection{Automatisation de lancement}

Le projet respecte l'esprit ``lancement simple'' demande dans le sujet:
\begin{itemize}
  \item scripts \texttt{run/validate/stop} pour Windows, Linux, macOS,
  \item verification prerequis (docker, kubectl, cluster),
  \item construction images + apply kustomize + attentes de readiness,
  \item ouverture de tunnels locaux et script de validation final.
\end{itemize}

\section{CI/CD et qualite de livraison}

Le workflow GitHub Actions couvre les etapes essentielles:
\begin{enumerate}
  \item \textbf{backend-ci}: installation deps Python, \texttt{manage.py check --deploy}, verification syntaxe,
  \item \textbf{frontend-ci}: \texttt{npm ci} puis build,
  \item \textbf{k8s-validate}: rendu kustomize et verification non-vide,
  \item \textbf{docker-build}: build des images backend/frontend,
  \item \textbf{publish-images} (manuel): push Docker Hub.
\end{enumerate}

\textbf{Interpretation professionnelle:} le pipeline n'est pas cosmetique; il detecte des regressions de configuration,
assure la buildabilite et valide la chaine de deploiement.

\section{Securite, limites et ameliorations}

\subsection{Points securite deja presents}
\begin{itemize}
  \item Separation ConfigMap/Secrets,
  \item Politique reseau restrictive (deny-all puis exceptions minimales),
  \item Conteneur backend execute en non-root,
  \item JWT court + refresh token,
  \item verification email cote inscription.
\end{itemize}

\subsection{Limites identifiees}
\begin{itemize}
  \item Pas de monitoring centralise (Prometheus/Grafana) dans le scope courant,
  \item Pas de tests d'integration end-to-end automatiques,
  \item Publication Docker en mode manuel (\texttt{workflow\_dispatch}) et non sur chaque push,
  \item Configuration locale orientee demo (pas de TLS de production actif dans le cluster local).
\end{itemize}

\subsection{Plan d'amelioration credible}
\begin{enumerate}
  \item Ajouter tests API et tests front automatises au pipeline,
  \item Integrer observabilite (metrics + dashboards),
  \item Ajouter deploiement automatise vers un cluster cible (staging),
  \item Renforcer politique de secrets (external secret manager),
  \item Mettre en place une strategie de rollback versionnee.
\end{enumerate}

\section{Conclusion}

Le projet repond au cahier des charges de maniere complete et argumentee:
\begin{itemize}
  \item architecture coherente et justifiee,
  \item fonctionnalites minimales implementees + options valorisees,
  \item conteneurisation Docker et orchestration Kubernetes operationnelles,
  \item documentation, scripts et pipeline permettant une evaluation reproductible.
\end{itemize}

Au-dela de la simple ``application qui tourne'', le livrable demontre une approche d'ingenierie logicielle
structuree: choix explicites, compromis assumes, automatisation, et maitrise de la chaine complete du developpement au deploiement.

\vspace{0.5cm}
\noindent\textbf{Annexe pratique (commande demo rapide):}
\begin{itemize}
  \item Windows: \texttt{run-project.bat} puis \texttt{validate-project.bat}
  \item Linux: \texttt{./scripts/linux/run-project.sh} puis \texttt{./scripts/linux/validate-project.sh}
  \item macOS: \texttt{./scripts/macos/run-project.sh} puis \texttt{./scripts/macos/validate-project.sh}
\end{itemize}

\end{document}
